\section{Experiments}
\label{sec:experiment}
We run Algorithm~\ref{alg:apriori} on a synthetic transaction database and vary the
minimum support $\sigma$. Figure~\ref{fig:runtime} shows that the runtime drops
sharply as $\sigma$ increases, because the pruning of Lemma~\ref{lem:apriori}
removes more candidates early. For comparison we also plot the FP-growth baseline
\cite{Han2000}.

\begin{figure}[ht]
\centering
\includegraphics[width=0.7\linewidth]{runtime_apriori.png}
\caption{Mining runtime versus the minimum support threshold.}
\label{fig:runtime}
\end{figure}

Table~\ref{tab:results} reports, for each support threshold, the number of frequent
itemsets and association rules produced by Algorithm~\ref{alg:rules}.

\begin{table}[ht]
\centering
\caption{Frequent itemsets and rules found at each support threshold.}
\label{tab:results}
\begin{tabular}{rrr}
\toprule
\textbf{Support (\%)} & \textbf{Frequent itemsets} & \textbf{Rules} \\
\midrule
1  & 312 & 1045 \\
2  & 148 & 392 \\
5  & 47  & 86 \\
10 & 13  & 18 \\
\bottomrule
\end{tabular}
\end{table}

As summarised in Figure~\ref{fig:runtime} and Table~\ref{tab:results}, a higher
support threshold yields fewer patterns and faster mining, which matches the
pruning behaviour analysed in Section~\ref{sec:method}.
